\indent ಮೇಲಿನ ವಾಕ್ಯಗಳಲ್ಲಿ ಕಂಡುಬರುವ ಕ್ರಿಯಾಪದಗಳಿಂದ ಕ್ರಿಯೆ ಈಗ ನಡೆಯುತ್ತಿರುವುದು ಎಂದು ತಿಳಿದು ಬರುವ ಕಾರಣ, ಈ ಕ್ರಿಯಾಪದಗಳಿಗೆ ವರ್ತಮಾನ ಕಾಲದ ಕ್ರಿಯಾಪದಗಳೆಂದು ಹೆಸರು. ಆದುದರಿಂದ ಮೇಲಿನ ಕ್ರಿಯಾಪದಗಳಲ್ಲಿರುವ ಆನೆ, ಆಳೆ, ಅದೆ, ಈಯೆ, ಏನೆ, ಆರೆ, ಅವೆ, ಈರಿ, ಏವೆ—ಎಂಬುವು ವರ್ತಮಾನ ಕಾಲದ ಆಖ್ಯಾತ ಪ್ರತ್ಯಯಗಳು\par

\vspace{4pt plus 4pt minus 1pt}\noindent{\kbvsecfont\fontsize{11pt}{14.5pt}\selectfont\color{kbvaccentdark} (2)\hspace{0.45em}ಭೂತ ಕಾಲದ ಕ್ರಿಯಾಪದ ರೂಪಗಳು :}\hspace{0.5em}ಗತಿಸಿಹೋದ ಕಾಲದಲ್ಲಿ ಕ್ರಿಯೆ ನಡೆಯಿತೆಂದು ತಿಳಿಸುವ ಕ್ರಿಯಾಪದವು. ಭೂತಕಾಲದ ಕ್ರಿಯಾಪದವೆನಿಸುವುದು.\par

\indent ಇಲ್ಲಿ ಕೊಟ್ಟಿರುವ ವಾಕ್ಯಗಳಲ್ಲಿ ಕಂಡುಬರುವ ಕ್ರಿಯಾಪದಗಳನ್ನು ಗಮನಿಸಿರಿ.\par

\indent (ಪ್ರ. ಪು. ಏ.)... ಅವನು ಓಡಿದನು -- ಓಡು + ದ + ಅನು\par

\indent (ಪ್ರ. ಸ್ತ್ರೀ. ಏ.)... ಅವಳು ಓಡಿದಳು -- ಓಡು + ದ + ಅಳು\par

\indent (ಪ್ರ. ನ. ಏ)... ಅದು ಓಡಿತು... ಓಡು + ದ + ಇತು\par

\indent (ಮ. ವಾ. ಏ)... ನೀನು ಓಡಿದೆ... ಓಡು + ದ + ಎ\par

\indent (ಉ. ವಾ. ಏ.)... ನಾನು ಓಡಿದೆನು... ಓಡು + ದ + ಎನು\par

\indent (ಪ್ರ. ವಾ. ಬ)... ಅವರು ಓಡಿದರು... ಓಡು + ದ + ಅರು\par

\indent (ಪ್ರ. ನ. ಬ)... ಅವು ಓಡಿದವು... ಓಡು + ದ + ಅವು\par

\indent (ಮ. ವಾ. ಬ.)... ನೀವು ಓಡಿದಿರಿ... ಓಡು + ದ + ಇರಿ\par

\indent (ಉ. ವಾ. ಬ)... ನಾವು ಓಡಿದೆವು... ಓಡು + ದ + ಎವು\par

\indent ಇಲ್ಲಿ ಒಂಬತ್ತು ಕ್ರಿಯಾಪದಗಳಿವೆ. ಈ ಕ್ರಿಯಾಪದಗಳಲ್ಲಿ “ಓಡು” ಧಾತುವಿನ ಮುಂದೆ ಎಲ್ಲ ಕಡೆಗೂ "ದ" ಎಂಬ ಪ್ರತ್ಯಯ ಬಂದಿದೆ ; ಪ್ರತ್ಯಯದ ಮುಂದೆ ಏಕವಚನದಲ್ಲಿ ಅನು, ಅಳು, ಇತು, ಎ, ಎನು, ಎಂಬ ಆಖ್ಯಾತ ಪ್ರತ್ಯಯಗಳು ಬಹುವಚನದಲ್ಲಿ ಅರು, ಅವು, ಇರಿ, ಎವು, ಎಂಬ ಆಖ್ಯಾತ ಪ್ರತ್ಯಯಗಳು ಬಂದಿವೆ. ಮೂಲ ಧಾತುವಿಗೂ ಆಖ್ಯಾತ ಪ್ರತ್ಯಯಕ್ಕೂ ನಡುವೆ ಬಂದಿರುವ "ದ" ಎಂಬುದು ಕಾಲಸೂಚಕ ಪ್ರತ್ಯಯ. ಇದು ಭೂತಕಾಲವನ್ನು ಸೂಚಿಸುವುದು. ಅಂದರೆ, ಕ್ರಿಯಾಪದವು ಭೂತಕಾಲದಲ್ಲಿರುವುದೆಂಬುದನ್ನು ತಿಳಿಸುವುದು.\par
